\setcounter{section}{0}
\Section{Счетность множества рациональных чисел, несчетность множества
действительных (вещественных) чисел}

\begin{definition}
Рациональное число~--- это дробь вида \(\frac{p}{q}\), где \(p \in \Z\),
\(q \in N\)
\end{definition}

Аксиомы рациональных чисел:

\begin{properties}
\item Если \(a = b\) и \(b = c\), то \(a = c\),

Если \(a > b\) и \(b > c\), то \(a > c\).

\item \(a + b = b + a\);
\item \((a + b) + c = a + (b + c)\);
\item Существует число~\(0\) такое, что для любого~\(a\) выполняется \(a + 0 = a\)
\item Для любого числа~\(a\) найдется противоположное число~\(-a\) такое,
что \(a + (-a) = 0\);
\item \(a \cdot b = b \cdot a\);
\item \((a \cdot b) \cdot c = a \cdot (b \cdot c)\);
\item Существует число 1 такое, что для любого~\(a\) выполняется \(a \cdot 1 = a\);
\item Для любого~\(a \neq 0\) существует обратное число \(\dfrac{1}{a}\) такое,
что \(a \cdot \dfrac{1}{a} = 1\);
\item \((a + b) \cdot c = a \cdot c + b \cdot c\);
\item Если \(a > b\), то для любого~\(c\) выполняется \(a + c > b + c\);
\item Если \(a > b\), то для любого~\(c > 0\) выполняется \(a \cdot c > b \cdot c\);
\item Для любого~\(a > 0\) найдется \(n \in \mathbb{N}\) такое, что \(n > a\).
\end{properties}

\begin{definition}
Сечением множества \(\mathbb{Q}\) называется разбиение этого множества
на 2 непустых множества \(A_*, A^*\), где:

\textbf{(A)} \(A_*\)~--- нет наибольшего, \(A^*\)~--- есть наименьшее;

\textbf{(C)} \(A_*\)~--- нет наибольшего, \(A^*\)~--- нет наименьшего.

\end{definition}

\begin{definition}
Действительное число --- сечение вида \textbf{(А)} или \textbf{(С)}.
\end{definition}

\begin{definition}
\(A\) называется счётным, если можно установить биекцию между \(A\) и \(\N\).
\end{definition}

\begin{theorem}
\(\mathbb{Q}\) --- счётно.
\end{theorem}

\begin{proof}
Множество положительных несократимых дробей \(p/q\)

\(H = p + q\) --- высота дроби.

\[
\begin{array}{|c|c|cc|cc|cccc|c}
\hline
\text{\(\N\)}  & 1   & 2   & 3   & 4   & 5   & 6   & 7   & 8   & 9   & ... \\
\hline
\text{\(p/q\)} & 1/1 & 1/2 & 2/1 & 1/3 & 3/1 & 1/4 & 2/3 & 3/2 & 4/1 & ... \\
\hline
\text{\(H\)}   & 2 & \multicolumn{2}{c|}{3} & \multicolumn{2}{c|}{4} &
    \multicolumn{4}{c|}{5} & ... \\
\hline
\end{array}
\]

\[
a_1, a_2, \ldots, a_n, \ldots \Leftrightarrow  0,\ a_1,\ -a_1,\ a_2,\ -a_2,\
\ldots
\Rightarrow \mathbb{Q}\text{~--- счётно.}
\]
\end{proof}

\begin{theorem}
Если \(A\) счётно, то и подмножество \(B \subset A\) счётно
(\(B\) не является конечным).
\end{theorem}

\begin{proof}
\[
\{a_1, a_2, a_3, a_4, \ldots, a_n, \ldots\} = A
\]

\[
\{a_{n_1}, a_{n_2}, a_{n_3}, a_{n_4}, \ldots, a_{n_k}, \ldots\} = B
\]

\[
a_n \leftrightarrow h
\]
\end{proof}

\begin{theorem}
\(\mathbb{R}\) --- несчётно.
\end{theorem}

\begin{proof}
Достаточно показать, что интервал \((0, 1)\)~---~несчётён.

Доказываем от противного: пусть все точки \((0, 1)\) пронумерованы
\(\Rightarrow\) каждую точку можно представить как
\[
\alpha_n = 0{,}a_1^n a_2^n a_3^n\ldots a_m^n\ldots
\]

Рассмотрим точку
\[
\gamma = 0{,}c_1 c_2 \ldots c_k \ldots, \qquad c_k \neq a_k^n
\]

\(\gamma \neq \alpha_n\ \forall n \Rightarrow (0, 1)\) --- несчётен.

Если бы \(\mathbb{R}\) было счётно, то интервал \((0, 1)\) был бы счётным
--- противоречие \(\Rightarrow \mathbb{R}\)~---~несчётно.
\end{proof}

\Section{Теорема о существовании точной верхней (нижней) грани множества}

\noindent\textbf{Важные определения:}

\begin{enumerate}
\item
\[
[\alpha_0 = \sup X] \eqdef
\begin{cases}
\forall x \in X \to \alpha_0 \ge x,\\
\forall \alpha' < \alpha_0 \ \exists x' \in X: x' > \alpha'.
\end{cases}
\]

\[
[\alpha_0 = \sup X] \eqdef  
[\forall x \in X \to \alpha_0 \ge x] \ \&\
[\forall \varepsilon > 0 \ \exists x_\varepsilon \in X: x_\varepsilon >
\alpha_0 - \varepsilon].
\]

\item 
\[
[\beta_0 = \inf X] \eqdef
\begin{cases}
\forall x \in X \to \beta_0 \le x,\\
\forall \beta' > \beta_0 \ \exists x' \in X: \beta' > x'.
\end{cases}
\]

\[
[\beta_0 = \inf X] \eqdef
\begin{cases}
\forall x \in X \to \beta_0 \le x,\\
\forall \varepsilon > 0 \ \exists x_\varepsilon \in X: x_\varepsilon < \beta_0
+ \varepsilon.
\end{cases}
\]

\end{enumerate}

\begin{theorem} [О существовании точной верхней (нижней) грани]
\[
[X \subset \mathbb{R}: \left[X \neq \varnothing \right] \& \left[X
\text{ огр. сверху}\right]] \Rightarrow
[\exists \alpha_0 \in \mathbb{R}: \alpha_0 = \sup X].
\]
\[
[X \subset \mathbb{R}: \left[X \neq \varnothing \right] \& \left[X
\text{ огр. снизу}\right]]
\Rightarrow
[\exists \beta_0 \in \mathbb{R}: \beta_0 = \inf X].
\]
\end{theorem}

\begin{proof}
Доказательство приведём для точной верхней грани (для нижней аналогично):
\begin{enumerate}
\item \(X\) имеет наибольший элемент \(x_0\);
\item \(X\) не имеет наибольшего элемента.
\end{enumerate}

Тогда:

\begin{enumerate}
\item \(\forall x \in X \ x \le x_0 \Rightarrow \forall \alpha \ge x_0
\Leftrightarrow x_0\) --- наименьшая из верхних граней;
\item Сечение \(\mathbb{R}\): \(A^*\) --- мн-во верхних граней,
\(A_* \neq \varnothing\) так как \(X\) ограничено. \(X \subset A_*\). При этом в
\(A_*\) нет наибольшего элемента. По теореме Дедекинда: \(\exists \alpha_0 \in
A^* \Rightarrow \alpha_0 = \sup X\).
\end{enumerate}
\end{proof}

\Section{Теорема об отделимости двух множеств действительных чисел}

\begin{theorem}
Пусть \(A \neq \varnothing\), \(B \neq \varnothing\), \(A, B \subset \mathbb{R}\),
и для любых \(a \in A\), \(b \in B\) выполнено \(b \ge a\). Тогда:
\[
\exists \alpha_0 = \sup A,\qquad \exists \beta_0 = \inf B,\qquad
\beta_0 \ge \alpha_0.
\]
\end{theorem}

\begin{proof}
Так как \(\beta_0 = \inf B\), то для любого \(b \in B\) число \(b\) является
верхней гранью множества \(A\) (по определению: \(\forall a \in A\ a \le b\)).

Поскольку \(A \neq \varnothing\) и \(A\) ограничено сверху, по теореме о
существовании точной верхней грани
\[
\exists \alpha_0 = \sup A.
\]
Аналогично, \(B \neq \varnothing\) и \(B\) ограничено снизу (например, любым
\(a \in A\)), поэтому
\[
\exists \beta_0 = \inf B.
\]
Далее, для любого \(b \in B\) имеем \(b \ge \alpha_0\), так как \(\alpha_0\) ---
наименьшая верхняя грань \(A\), а каждое \(b\) --- верхняя грань \(A\). Значит,
\(\alpha_0\) --- нижняя грань \(B\). Следовательно,
\[
\beta_0 = \inf B \ge \alpha_0.
\]
\end{proof}

\Section{Единственность предела сходящейся последовательности. Ограниченность
сходящейся последовательности}


\begin{definition}
Последовательность называется сходящейся, если у неё есть конечный предел.
\end{definition}

\begin{definition}
\[
\left[\lim_{n\to\infty} x_n = \alpha\right] \eqdef
\bigl[\forall \epsilon>0, \exists N = N(\epsilon)\colon \forall n\ge N
\to |x_n-\alpha|<\epsilon\bigr].
\]
\end{definition}

\begin{theorem}[единственность предела]
Если последовательность \(x_n\) сходится, то она имеет единственный предел.
\end{theorem}

\begin{proof}
Пусть существуют два предела \(\alpha\) и \(\beta\). Тогда
\[
x_n = \alpha + \alpha_n,\qquad x_n = \beta + \beta_n,
\]
где \(\alpha_n,\beta_n\) --- бесконечно малые. Тогда
\[
\gamma = \alpha-\beta
= x_n-\alpha_n - x_n + \beta_n
= \beta_n-\alpha_n
= \gamma_n.
\]
Так как \(\gamma=\gamma_n\), то \(\gamma=0\), значит \(\alpha=\beta\).
Противоречие с \(a\ne\beta\). Следовательно, предел единственен.
\end{proof}

\begin{theorem}[ограниченность сходящейся последовательности]
Если последовательность \(\{x_n\}\) сходится, то она ограничена.
\end{theorem}

\begin{proof}
Пусть \(\exists \alpha\in\mathbb R\) такое, что
\(\forall n\ x_n=\alpha+\alpha_n\), где
\(\alpha_n\) --- бесконечно малая. Тогда:
\[
|x_n| = |\alpha+\alpha_n| \le |\alpha| + |\alpha_n|.
\]
Так как \(\alpha_n\) ограничена, \(\exists C>0:\ |\alpha_n|\le C\).
Следовательно,
\[
|x_n| \le |\alpha| + C,
\]
то есть \(\{x_n\}\) ограничена.
\end{proof}

Обратное неверно: последовательность \((-1)^n\) ограничена, но не сходится.

\Section{Бесконечно малые последовательности и их свойства}

\textbf{Обозначения:} \(\text{б.м.}\) --- бесконечно малая,
\(\text{б.б.}\) --- бесконечно большая.

\begin{definition}
\[
\left[\left\{x_n\right\}~-\text{б.м.}\right] \eqdef
\left[\forall\epsilon > 0\ \exists N = N(\epsilon) \colon \forall n
\geq N\colon |x_n| < \epsilon\right]
\]
\end{definition}

\textbf{Свойства:}

\begin{enumerate}
\item Если \(\{x_n\}\) и \(\{y_n\}\) --- б.м., то \(\{z_n = x_n + y_n\}\) --- б.м.

\begin{proof}
Так как \(x_n, y_n\) --- б.м., то
\[
\forall \varepsilon>0\ \exists N_1\in\mathbb{N}:\ \forall n\ge N_1\
|x_n|<\frac{\varepsilon}{2},
\]
\[
\forall \varepsilon>0\ \exists N_2=N_2(\varepsilon):\ \forall n\ge N_2\
|y_n|<\frac{\varepsilon}{2}.
\]
Положим \(N=\max(N_1,N_2)\). Тогда \(\forall n\ge N\)
\[
|z_n|=|x_n+y_n|\le |x_n|+|y_n|<\frac{\varepsilon}{2} +
\frac{\varepsilon}{2}=\varepsilon.
\]
Следовательно, \(z_n\) --- б.м.

\end{proof}

\item Если \(\{x_n\}\) --- б.м., то \(\{x_n\}\) ограничена.

\begin{proof} Возьмём \(\varepsilon=1\). Тогда
\[
\exists N=N(1):\ \forall n\ge N\ |x_n|<1.
\]
Положим
\[
C=\max(|x_1|,|x_2|,\ldots,|x_{N-1}|,1).
\]
Тогда \(\forall n\in\mathbb{N}\ |x_n|\le C\). Значит, \(\{x_n\}\) ограничена.

\end{proof}

\item Если \(\{x_n\}\) --- б.м., а \(\{y_n\}\) ограничена, то \(\{z_n=y_n x_n\}\)
--- б.м.

\begin{proof} Так как \(\{y_n\}\) ограничена, \(\exists C>0:\ \forall
n \in\mathbb{N}\ |y_n|\le C\). Для любого \(\varepsilon>0\) найдётся
\(N=N(\varepsilon)\) такое, что \(\forall n\ge N\ |x_n|<\frac{\varepsilon}{C}\).
Тогда:
\[
|z_n|=|x_n|\cdot|y_n|\le C\cdot\frac{\varepsilon}{C}<\varepsilon.
\]
Следовательно, \(z_n\) --- б.м.

\end{proof}

\item Если \(\{z_n\}\) --- б.м. и \(z_n=r\) при всех \(n\ge n_0\), то \(r=0\).

\begin{proof} Пусть \(r\ne 0\). Возьмём \(\varepsilon=\frac{|r|}{2}\). Тогда
\[
\exists N=N(\varepsilon)\ge n_0:\ \forall n\ge N\ |z_n|<\varepsilon.
\]
Но \(z_n=r\), поэтому \(|r|<\frac{|r|}{2}\) --- противоречие. Значит, \(r=0\).
\end{proof}

\item
\begin{enumerate}
\item[(5.1)] \(\displaystyle \frac{1}{\text{б.б.}}=\text{б.м.}\)

Пусть \(\{y_n\}\) --- б.б. и \(y_n\ne 0\) при \(n\ge n_0\). Тогда
\[
\forall \varepsilon>0\ \exists N=N(\varepsilon)\ge n_0:\ \forall n\ge N\
|y_n|>\frac{1}{\varepsilon}.
\]
Следовательно,
\[
|z_n|=\frac{1}{|y_n|}<\varepsilon,
\]
то есть \(z_n=\frac{1}{y_n}\) --- б.м.

\item[(5.2)] \(\displaystyle \frac{1}{\text{б.м.}}=\text{б.б.}\) --- аналогично.
\end{enumerate}

\item Следствие из пунктов 3 и 5:
\[
\frac{x_n}{y_n}=z_n-\text{б.м.},\qquad y_n-\text{б.б.}
\]
(то есть частное б.м. и б.б. есть б.м.)

\end{enumerate}

\Section{Арифметические операции со сходящимися последовательностями}

\begin{definition}
Число \(a\) называется пределом последовательности \(\{x_n\}\), если
\[
\lim_{n\to\infty} x_n = \alpha
\;\overset{\text{def}}{\Longleftrightarrow}\;
\forall \varepsilon>0\ \exists N=N(\varepsilon)\in\mathbb N\
\forall n\ge N:\ |x_n-\alpha|<\epsilon.
\]
\end{definition}

\begin{theorem}
Пусть последовательности \(\{x_n\}\) и \(\{y_n\}\) сходятся, причём
\[
\lim_{n\to\infty} x_n = \alpha,\qquad
\lim_{n\to\infty} y_n = \beta.
\]
Тогда:
\begin{enumerate}
\item \(\displaystyle \lim_{n\to\infty}(x_n\pm y_n)=\alpha\pm\beta\);
\begin{proof}
Если \(x_n = \alpha + \alpha_n\) и \(y_n = \beta + \beta_n\), то
\(z_n = x_n + y_n = (\alpha + \beta) + (\alpha_n + \beta_n)
\xrightarrow{n \rightarrow \inf} \alpha + \beta\)
\end{proof}

\item \(\displaystyle \lim_{n\to\infty}(x_n y_n)=\alpha\beta\);
\begin{proof}
Если \(x_n = \alpha + \alpha_n\) и \(y_n = \beta + \beta_n\), то:
\[
z_n = x_n \cdot y_n = (\alpha \cdot \beta) + (\alpha \cdot \beta_n +
\beta \cdot \alpha_n +\alpha_n \cdot \beta_n) \xrightarrow{n \rightarrow \inf}
\alpha + \beta
\]
\end{proof}

\item если \(\beta\ne 0\), то
\[
\lim_{n\to\infty}\frac{x_n}{y_n}=\frac{\alpha}{\beta}.
\]
\end{enumerate}
\end{theorem}

\begin{proof} Так как \(x_n=\alpha+\alpha_n\) и \(y_n=\beta+\beta_n\), то:
\[
\frac{x_n}{y_n}-\frac{\alpha}{\beta}
=
\frac{\beta x_n-\alpha y_n}{\beta y_n}
=
\frac{\beta(\alpha+\alpha_n)-\alpha(\beta+\beta_n)}{\beta y_n}
=
\frac{\beta\alpha_n-\alpha\beta_n}{\beta y_n}
=
\frac{1}{y_n}\left(\alpha_n-\frac{\alpha}{\beta}\beta_n\right)
=
\frac{1}{y_n}\delta_n,
\]
где:
\[
\delta_n=\alpha_n-\frac{\alpha}{\beta}\beta_n
\]

Поскольку \(y_n\to\beta\ne 0\), последовательность \(\frac{1}{y_n}\)
ограничена при достаточно больших \(n\). Поэтому произведение ограниченной
последовательности на бесконечно малую \(\delta_n\) есть бесконечно малая.
Следовательно,

\[
\frac{x_n}{y_n}-\frac{\alpha}{\beta}\to 0,
\]
то есть:
\[
\lim_{n\to\infty}\frac{x_n}{y_n}=\frac{\alpha}{\beta}.
\]
\end{proof}

\Section{Свойства пределов, связанные с неравенствами}

\begin{enumerate}
    \item Если \(x_n \to \alpha\) и \(\exists N: \forall n \ge N \implies x_n \ge \beta\), то и \(\alpha \ge \beta\).

\begin{proof}
\[\varepsilon = \beta - \alpha > 0\].
\end{proof}

\noindent\textbf{Следствия:}

\begin{enumerate}

\item Если \(\exists N: \forall n \ge N \implies x_n \ge y_n\), то и
\(
\lim x_n \ge \lim y_n.
\)

\begin{proof}
Рассмотрим \(z_n = x_n - y_n \ge 0\).
\end{proof}

\item Если \(\forall n \ge N \implies \alpha \le x_n \le \beta\), то
\( \alpha \le \lim x_n \le \beta. \)

\end{enumerate}

\item \([x_n \to \alpha] \ \&\ [z_n \to \alpha] \ \&\ [\exists N_0: \forall n \ge N_0 \implies x_n \le y_n \le z_n] \implies [y_n \to \alpha]\).

\begin{proof}
\[
x_n - \alpha \le y_n - \alpha \le z_n - \alpha
\implies
|y_n - \alpha| \le \max(|x_n - \alpha|, |z_n - \alpha|).
\]
\[
\forall \varepsilon > 0 \ \exists n_1: \forall n \ge n_1 \implies |x_n - \alpha| < \varepsilon.
\]
\[
\forall \varepsilon > 0 \ \exists n_2: \forall n \ge n_2 \implies |z_n - \alpha| < \varepsilon.
\]
Выберем \(N = \max(N_0, n_1, n_2)\). Тогда для всех \(n \ge N\) выполнено
\[
|y_n - \alpha| \le \max(|x_n - \alpha|, |z_n - \alpha|) < \varepsilon.
\]
Следовательно, \(y_n \to \alpha\).

\end{proof}

\end{enumerate}

\Section{Теорема о пределе ограниченной монотонной последовательности}

\begin{definition}
Последовательность называется неубывающей (невозрастающей), если
\[
\forall m, n \in \N,\ m > n \mapsto x_m \geq x_n
\]
\[
\forall m, n \in \N,\ m > n \mapsto x_m \leq x_n
\]

\end{definition}

\begin{theorem}
Если последовательность \(\{x_n\}\) монотонна (строго монотонна) и ограничена,
то существует \(\lim x_n\).

При этом если \(\{x_n\}\) возрастает (не убывает), то
\[
\lim_{n\to\infty} x_n = \sup x_n,
\]
иначе
\[
\lim_{n\to\infty} x_n = \inf x_n.
\]
\end{theorem}

\begin{proof}
Пусть \(\{x_n\}\) возрастает и ограничена сверху. Тогда существует
\[
\sup x_n = \alpha.
\]
По определению точной верхней грани
\[
\forall \varepsilon>0\ \exists N=N(\varepsilon):\ x_N > \alpha-\varepsilon.
\]
Так как последовательность возрастает, то для всех \(n\ge N\)
\[
x_N \le x_n \le \alpha.
\]
Следовательно,
\[
0 \le \alpha -x_n \le \alpha - x_N < \varepsilon.
\]
Значит,
\[
|x_n-\alpha|<\varepsilon \quad \text{при } n\ge N,
\]
то есть
\[
\lim_{n\to\infty} x_n = \alpha = \sup x_n.
\]

Для других случаев доказательство аналогично.
\end{proof}

\noindent\textbf{Замечание.} Так как любая сходящаяся последовательность
ограничена, то монотонная сходящаяся последовательность ограничена.

\Section{Теорема Кантора о вложенных отрезках}

\begin{theorem}[Кантора]
Пусть дана последовательность вложенных отрезков
\[
[a_{n+1}, b_{n+1}] \subset [a_n, b_n] \quad \forall n,
\]
причём
\[
\lim_{n\to\infty}(b_n-a_n)=0.
\]
Тогда существует единственная точка \(\alpha\), принадлежащая всем отрезкам:
\[
\exists! \alpha:\quad \alpha \in [a_n,b_n]\ \ \forall n.
\]
\end{theorem}

\begin{proof}
Так как \([a_{n+1},b_{n+1}]\subset [a_n,b_n]\), то
\[
a_n \le a_{n+1} \le b_{n+1} \le b_n.
\]
Следовательно, \(\{a_n\}\) не убывает и ограничена сверху, например, числом
\(b_1\), а \(\{b_n\}\) не возрастает и ограничена снизу, например, числом
\(a_1\). Поэтому существуют пределы
\[
\lim_{n\to\infty} a_n = \sup a_n = \alpha,\qquad
\lim_{n\to\infty} b_n = \inf b_n = \beta.
\]
Тогда
\[
\lim_{n\to\infty}(b_n-a_n)=\beta-\alpha.
\]
По условию этот предел равен нулю, значит,
\[
\beta-\alpha=0 \quad\Rightarrow\quad \beta=\alpha.
\]
Обозначим общее значение через \(\gamma\). Тогда
\[
a_n \le \gamma \le b_n \quad \forall n,
\]
то есть \(\gamma\in[a_n,b_n]\) для всех \(n\).

Докажем единственность. Пусть существует другая точка \(\gamma'\), принадлежащая
всем отрезкам, причём \(\gamma'>\gamma\). Тогда для всех \(n\)
\[
\gamma'-\gamma \le b_n-a_n.
\]
Но \(b_n-a_n\to 0\), значит, \(b_n-a_n\) — бесконечно малая. Следовательно,
\(\gamma'-\gamma=0\), то есть \(\gamma'=\gamma\). Противоречие. Значит, общая
точка единственна.
\end{proof}

\Section{Подпоследовательности и частичные пределы. Критерий частичного предела}

\begin{definition}
Пусть дана последовательность \(\{x_n\}\). Последовательность \(\{x_{n_k}\}\)
называется \emph{подпоследовательностью} последовательности \(\{x_n\}\), если
\[
n_1 < n_2 < \ldots < n_k < \ldots
\]
— строго возрастающая последовательность натуральных чисел. При этом \(n_k \ge k\).
\end{definition}

\textbf{Свойства:}

\begin{enumerate}
    \item
        \[
        \left[\exists \lim_{n\to\infty}x_n = \alpha\right] \Rightarrow
        \left[\forall x_{n_k} \mapsto \lim_{k\to\infty}x_{n_k} = \alpha\right]
        \]
    \item
        \[
        \left[\forall x_{n_k} \mapsto \lim_{k\to\infty}x_{n_k}
        \text{~--- сходится}\right] \Rightarrow
        \left[\forall x_{n_k} \mapsto \lim_{k\to\infty}x_{n_k} = \alpha \right]
        \]
    \item
        \[
        \left[\left\{x_n\right\}\text{~--- б.б.}\right]\Rightarrow
        \left[\forall\left\{x_{n_k}\right\}\text{~--- б.б.}\right]
        \]
    \item
        \[
        \left[\lim_{n\to\infty}x_n = \alpha\right]\Rightarrow
        \left[\exists\ \text{монот. } x_{n_k}\colon \lim_{k\to\infty}x_{n_k}
        = \alpha\right]
        \]
\end{enumerate}

\begin{definition}
Число \(\alpha\) называется \emph{частичным пределом} последовательности
\(\{x_n\}\), если существует подпоследовательность \(\{x_{n_k}\}\), сходящаяся к
\(\alpha\).
\end{definition}

\begin{definition}
    \(\alpha\)~--- предельная точка, если \(\forall U_\epsilon(\alpha)\)
\end{definition}

\begin{theorem}[критерий частичного предела]
Число \(a\) является частичным пределом последовательности \(\{x_n\}\) тогда и
только тогда, когда для любого \(\varepsilon>0\) и любого \(N\in\mathbb N\)
найдётся \(n>N\) такое, что
\[
|x_n - a| < \varepsilon.
\]
Иными словами, множество
\[
\{n\in\mathbb N: |x_n-a|<\varepsilon\}
\]
бесконечно для каждого \(\varepsilon>0\).
\end{theorem}

\begin{proof}
\emph{Необходимость.} Пусть \(a\) — частичный предел. Тогда существует
подпоследовательность \(x_{n_k}\to a\). Значит, для любого \(\varepsilon>0\)
найдётся \(K\) такое, что при \(k\ge K\) выполнено
\[
|x_{n_k}-a|<\varepsilon.
\]
Так как \(n_k\to\infty\), для любого \(N\) можно выбрать \(k\) настолько большим,
что \(n_k>N\). Тогда \(n=n_k>N\) и \(|x_n-a|<\varepsilon\).

\emph{Достаточность.} Пусть условие выполнено. Построим подпоследовательность,
сходящуюся к \(a\). Возьмём \(\varepsilon=1\) и \(N=0\). Найдётся \(n_1>0\)
с \(|x_{n_1}-a|<1\). Затем возьмём \(\varepsilon=1/2\) и \(N=n_1\). Найдётся
\(n_2>n_1\) с \(|x_{n_2}-a|<1/2\). Продолжая, получим
\[
n_1<n_2<\ldots
\]
и
\[
|x_{n_k}-a|<\frac{1}{k}.
\]
Тогда \(x_{n_k}\to a\), то есть \(a\) — частичный предел.
\end{proof}

\Section{Теорема Больцано–Вейерштрасса}

\begin{theorem}[Больцано--Вейерштрасса]
Из любой ограниченной последовательности можно выделить сходящуюся
подпоследовательность.
\end{theorem}

\begin{proof}
Пусть последовательность \(\{x_n\}\) ограничена:
\[
\exists a,b\in\mathbb R:\quad \forall n\in\mathbb N \quad a\le x_n\le b.
\]
Рассмотрим отрезок \([a,b]\). Разделим его пополам. По крайней мере в одной из
половин содержится бесконечно много членов последовательности. Выберем эту
половину и обозначим её \([a_1,b_1]\). Затем разделим \([a_1,b_1]\) пополам и
снова выберем ту половину, которая содержит бесконечно много членов
последовательности; обозначим её \([a_2,b_2]\). Продолжая этот процесс, получим
последовательность вложенных отрезков
\[
[a,b]\supset [a_1,b_1]\supset [a_2,b_2]\supset \ldots \supset [a_n,b_n]\supset
\ldots,
\]
причём
\[
b_n-a_n=\frac{b-a}{2^n}\to 0 \quad \text{при } n\to\infty.
\]
По теореме о вложенных отрезках существует единственная точка \(c\),
принадлежащая всем отрезкам:
\[
\exists! c:\quad c\in[a_n,b_n]\quad \forall n\in\mathbb N.
\]

Построим подпоследовательность. В каждом отрезке \([a_n,b_n]\) выберем член
последовательности \(x_{n_k}\) так, чтобы номера возрастали:
\[
\exists x_{n_1}\in[a_1,b_1],
\]
\[
\exists x_{n_2}\in[a_2,b_2],\quad n_2>n_1,
\]
\[
\ldots
\]
\[
\exists x_{n_k}\in[a_k,b_k],\quad n_k>n_{k-1}.
\]
Тогда для всех \(k\)
\[
a_k\le x_{n_k}\le b_k.
\]
Так как \(a_k\to c\) и \(b_k\to c\) (поскольку \(b_k-a_k\to 0\)), по теореме о
двух милиционерах получаем
\[
x_{n_k}\to c.
\]
Следовательно, из последовательности \(\{x_n\}\) выделена сходящаяся
подпоследовательность \(\{x_{n_k}\}\). Теорема доказана.
\end{proof}

\Section{Теорема о единственном частичном пределе}

\begin{theorem}
Пусть последовательность \(\{x_n\}\) ограничена и имеет единственный частичный
предел \(\alpha\). Тогда
\[
\lim_{n\to\infty} x_n = \alpha.
\]
\end{theorem}

\begin{proof}
Так как \(\{x_n\}\) ограничена, по теореме Больцано--Вейерштрасса из неё можно
выделить сходящуюся подпоследовательность:
\[
\exists \{x_{n_k}\}:\quad x_{n_k}\to \alpha.
\]
Значит, \(\alpha\) действительно является частичным пределом.

Докажем, что вся последовательность сходится к \(\alpha\). Предположим противное:
\[
x_n \not\to \alpha.
\]
Тогда существует \(\varepsilon>0\) такое, что вне \(\varepsilon\)-окрестности
точки \(\alpha\), то есть вне интервала
\[
(\alpha-\varepsilon,\ \alpha+\varepsilon),
\]
находится бесконечно много членов последовательности.

Так как \(\{x_n\}\) ограничена, все эти члены лежат в некотором отрезке
\([a,b]\). Вне \(\varepsilon\)-окрестности точки \(\alpha\) остаются два
промежутка:
\[
[a,\ \alpha-\varepsilon]
\qquad\text{и}\qquad
[\alpha+\varepsilon,\ b].
\]
По крайней мере в одном из них содержится бесконечно много членов
последовательности.

Пусть, например, бесконечно много членов лежит в \([\alpha+\varepsilon,\ b]\).
Тогда можно выбрать подпоследовательность
\[
x_{n_k}\in[\alpha+\varepsilon,\ b].
\]
Она ограничена, поэтому по теореме Больцано--Вейерштрасса из неё можно выделить
сходящуюся подпоследовательность:
\[
x_{n_{k_j}}\to \beta.
\]
При этом
\[
\beta\in[\alpha+\varepsilon,\ b],
\]
значит,
\[
\beta\neq \alpha.
\]
Но тогда \(\beta\) --- ещё один частичный предел последовательности \(\{x_n\}\),
что противоречит условию единственности частичного предела.

Аналогично рассматривается случай, когда бесконечно много членов лежит в \([a,\
\alpha-\varepsilon]\).

Полученное противоречие показывает, что предположение \(x_n \not\to \alpha\)
неверно. Следовательно,
\[
\lim_{n\to\infty} x_n = \alpha.
\]
\end{proof}

\Section{Критерий Коши сходимости числовой последовательности}

\begin{definition}
Последовательность \(\{x_n\}\) называется \emph{фундаментальной}
(или последовательностью Коши), если
\[
\forall \varepsilon>0\ \exists N=N(\varepsilon)\in\mathbb N:\ 
\forall n>N\ \forall p\in\mathbb N
\quad |x_{n+p}-x_n|<\varepsilon.
\]
\end{definition}

\begin{theorem}
Если \(\{x_n\}\) фундаментальна, то она ограничена.
\end{theorem}

\begin{proof}
Возьмём \(\varepsilon=1\). Тогда
\[
\exists n\in\mathbb N\ \forall p\in\mathbb N:\ |x_{n+p}-x_n|<1.
\]
Следовательно,
\[
|x_{n+p}|=|x_{n+p}-x_n+x_n|
\le |x_{n+p}-x_n|+|x_n|
<1+|x_n|.
\]
Положим
\[
C=\max\bigl(|x_1|,|x_2|,\ldots,|x_{n-1}|,\ |x_n|+1\bigr).
\]
Тогда для всех \(k\in\mathbb N\) выполнено \(|x_k|\le C\), то есть \(\{x_n\}\)
ограничена.
\end{proof}

\begin{theorem}[Критерий Коши]
Последовательность \(\{x_n\}\) сходится тогда и только тогда, когда она
фундаментальна:
\[
[\{x_n\}\text{ сходится}]
\iff
[\{x_n\}\text{ фундаментальна}].
\]
\end{theorem}

\begin{proof}
\emph{Необходимость.} Пусть \(x_n\to \alpha\). Тогда
\[
\forall \varepsilon>0\ \exists N=N(\varepsilon):\
\forall n>N\ |x_n-\alpha|<\frac{\varepsilon}{2}.
\]
Для любого \(p\in\mathbb N\) имеем \(n+p>N\), поэтому
\[
|x_{n+p}-x_n|
=
|x_{n+p}-\alpha-(x_n-\alpha)|
\le
|x_{n+p}-\alpha|+|x_n-\alpha|
<
\frac{\varepsilon}{2}+\frac{\varepsilon}{2}
=
\varepsilon.
\]
Значит, \(\{x_n\}\) фундаментальна.

\emph{Достаточность.} Пусть \(\{x_n\}\) фундаментальна. По прошлой теореме она
ограничена. Тогда по теореме Больцано--Вейерштрасса из неё можно выделить
сходящуюся подпоследовательность:
\[
x_{n_k}\to \alpha.
\]
Покажем, что вся последовательность сходится к \(\alpha\).
Для \(\varepsilon>0\) выберем \(N_1=N_1(\varepsilon)\) так, что
\[
\forall n,m\ge N_1:\ |x_n-x_m|<\frac{\varepsilon}{2}.
\]
Также выберем \(N_2=N_2(\varepsilon)\) так, что
\[
\forall n_k\ge N_2:\ |x_{n_k}-\alpha|<\frac{\varepsilon}{2}.
\]
Положим \(N=\max(N_1,N_2)\). Для любого \(n\ge N\) выберем \(k\) настолько
большим, что \(n_k\ge N\). Тогда
\[
|x_n-\alpha|
\le
|x_n-x_{n_k}|+|x_{n_k}-\alpha|
<
\frac{\varepsilon}{2}+\frac{\varepsilon}{2}
=
\varepsilon.
\]
Следовательно, \(x_n\to \alpha\).
\end{proof}

\Section{Определение предела функции в точке по Коши и по Гейне, их
эквивалентность}

\noindent\textbf{Определение (Гейне).}
\[
[\lim_{x\to\alpha} f(x)=\beta]
\eqdef
\bigl[\forall \{x_n\}:\ [\{x_n\}\to\alpha]\ \&\ [x_n\ne\alpha\ \forall n]
\ \Rightarrow\ f(x_n)\to\beta\bigr].
\]

\noindent Тогда
\[
y_n=f(x_n)\xrightarrow[n\to\infty]{}\beta.
\]

\noindent\textbf{Определение (Коши).}
\[
[\lim_{x\to\alpha} f(x)=\beta]
\eqdef
\bigl[\forall \varepsilon>0\ \exists \delta=\delta(\varepsilon)>0\
\forall x\in \dot U_\delta(\alpha)
\ \Rightarrow\ f(x)\in U_\varepsilon(\beta)\bigr].
\]

\begin{theorem}
\[
[\text{определение Коши}]
\iff
[\text{определение Гейне}].
\]
\end{theorem}

\begin{proof}
\noindent 1) Коши \(\Rightarrow\) Гейне.

Пусть \(\lim_{x\to\alpha} f(x)=\beta\) по Коши. Возьмём произвольную
последовательность \(\{x_n\}\) такую, что
\[
x_n\to\alpha,\qquad x_n\ne\alpha\ \forall n.
\]
Тогда
\[
\exists N=N(\delta)\ \forall n\ge N:\quad
0<|x_n-\alpha|<\delta
\ \Rightarrow\
|f(x_n)-\beta|<\varepsilon.
\]
Следовательно, \(f(x_n)\to\beta\), то есть выполнено определение Гейне. ч.т.д.

\medskip
\noindent 2) Гейне \(\Rightarrow\) Коши. От противного.

Пусть предел по Коши не равен \(\beta\). Тогда
\[
[\lim_{x\to\alpha} f(x)\ne\beta]
\eqdef
\bigl[\exists \varepsilon_0>0\ \forall \delta>0\
\exists x\in \dot U_\delta(\alpha):\
f(x)\notin U_{\varepsilon_0}(\beta)\bigr].
\]
Для каждого \(n\in\mathbb N\) положим \(\delta_n=\frac1n\). Тогда
\[
\exists x_n:\quad
0<|x_n-\alpha|<\delta_n
\ \Rightarrow\
|f(x_n)-\beta|\ge \varepsilon_0.
\]
При этом
\[
x_n\ne\alpha\ \forall n,\qquad x_n\to\alpha.
\]
Значит, \(\{x_n\}\) --- последовательность Гейне, и по определению Гейне должно
быть
\[
f(x_n)\to\beta.
\]
Но это противоречит неравенству \(|f(x_n)-\beta|\ge\varepsilon_0\). Полученное
противоречие доказывает, что
\[
\lim_{x\to\alpha} f(x)=\beta
\]
и по Коши.
\end{proof}

\Section{Критерий Коши существования предела функции}

\begin{theorem}[Критерий Коши]
Пусть функция \(f\) определена в проколотой окрестности точки \(\alpha\).
Тогда
\[
\bigl[\lim_{x\to \alpha} f(x)=\beta\in\mathbb R\bigr]
\iff
\bigl[\forall \varepsilon>0\ \exists \delta=\delta(\varepsilon)>0\
\forall x',x''\in \dot U_\delta(\alpha):\
|f(x')-f(x'')|<\varepsilon\bigr].
\]
\end{theorem}

\begin{proof}
\noindent\textbf{Необходимость.}
Пусть \(\lim_{x\to \alpha} f(x)=\beta\). Тогда
\[
\forall \varepsilon>0\ \exists \delta=\delta(\varepsilon)>0\
\forall x\in \dot U_\delta(\alpha):\
|f(x)-\beta|<\frac{\varepsilon}{2}.
\]
Тогда для любых \(x',x''\in \dot U_\delta(\alpha)\)
\[
|f(x')-f(x'')|
=
|(f(x')-\beta)-(f(x'')-\beta)|
\le
|f(x')-\beta|+|f(x'')-\beta|
<
\frac{\varepsilon}{2}+\frac{\varepsilon}{2}
=
\varepsilon.
\]

\medskip
\noindent\textbf{Достаточность.}
Пусть условие Коши выполнено. Покажем, что существует предел.

Возьмём произвольную последовательность \(\{x_n\}\) такую, что
\[
x_n\to \alpha,\qquad x_n\ne \alpha\ \ \forall n.
\]
По условию Коши для любого \(\varepsilon>0\) найдётся
\(\delta=\delta(\varepsilon)>0\).
Так как \(x_n\to \alpha\), существует \(N=N(\delta)\) такое, что
\[
\forall n\ge N:\quad 0<|x_n-\alpha|<\delta.
\]
Тогда для любых \(n,m\ge N\) имеем \(0<|x_n-\alpha|<\delta\) и
\(0<|x_m-\alpha|<\delta\), следовательно,
\[
|f(x_n)-f(x_m)|<\varepsilon.
\]
Значит, последовательность \(\{f(x_n)\}\) фундаментальна. По критерию Коши для
последовательностей она сходится. Обозначим её предел через \(\beta\).

Покажем, что \(\beta\) не зависит от выбора последовательности. Пусть
\(\{x_n'\}\) и \(\{x_n''\}\) --- две последовательности, сходящиеся к \(\alpha\),
причём
\[
f(x_n')\to \beta',\qquad f(x_n'')\to \beta''.
\]
Рассмотрим смешанную последовательность
\[
x_1',\,x_1'',\,x_2',\,x_2'',\,\ldots
\]
Она также сходится к \(\alpha\), и все её члены отличны от \(\alpha\). По
доказанному последовательность значений \(f\) на ней сходится. Но её
подпоследовательности \(f(x_n')\) и \(f(x_n'')\) сходятся соответственно к
\(\beta'\) и \(\beta''\). Следовательно, \(\beta'=\beta''\). Значит, предел один
и тот же для всех последовательностей, стремящихся к \(\alpha\).

По определению предела по Гейне
\[
\lim_{x\to \alpha} f(x)=\beta.
\]
Теорема доказана.
\end{proof}

\Section{Существование односторонних пределов у монотонных функций}

\begin{theorem}
Пусть функция \(f(x)\) ограничена на \((a,b)\) и монотонна на \((a,b)\).
Тогда для любой точки \(x_0\in(a,b)\) существуют односторонние пределы
\(f(x_0-0)\) и \(f(x_0+0)\), причём:
\begin{itemize}
\item если \(f\) неубывающая, то
\[
f(x_0-0)\le f(x_0)\le f(x_0+0);
\]
\item если \(f\) невозрастающая, то
\[
f(x_0-0)\ge f(x_0)\ge f(x_0+0).
\]
\end{itemize}
\end{theorem}

\begin{proof}
Пусть \(y=f(x)\) --- неубывающая функция. Зафиксируем точку \(x_0\in(a,b)\).
Для любого \(x_1\in(a,x_0)\) имеем
\[
f(x_1)\le f(x_0).
\]
Следовательно, функция \(f(x)\) ограничена на \((a,x_0)\), и существует
\[
\exists \sup_{(a,x_0)} f(x)=\alpha.
\]
При этом
\[
\alpha\le f(x_0),
\]
то есть для любого \(\varepsilon>0\) найдётся \(x_\varepsilon\in(a,x_0)\) такое,
что
\[
f(x_\varepsilon)>\alpha-\varepsilon.
\]
Положим
\[
\delta=x_0-x_\varepsilon.
\]
Тогда для всех \(x\in(x_0-\delta,x_0)\) выполнено
\[
\alpha-\varepsilon<f(x_\varepsilon)\le f(x)\le\alpha,
\]
значит,
\[
|f(x)-\alpha|<\varepsilon.
\]
Следовательно,
\[
\lim_{x\to x_0-0} f(x)=\alpha=f(x_0-0).
\]
Аналогично доказывается существование правого предела \(f(x_0+0)\).
Для невозрастающей функции доказательство проводится аналогично.
\end{proof}

\Section{Непрерывность функции в точке. Непрерывность сложной функции}

\begin{definition}
Функция \(f\) называется непрерывной в точке \(\alpha\), если
\[
[f \text{ непрерывна в т. } \alpha]
\eqdef
\bigl[\lim_{x\to\alpha} f(x)=f(\alpha)\bigr].
\]
\end{definition}

\begin{theorem}
Если функции \(f\) и \(g\) непрерывны в точке \(\alpha\), то функции
\(f+g\), \(fg\), а также \(\dfrac{f}{g}\) (при \(g(\alpha)\ne 0\)) непрерывны
в точке \(\alpha\).
\end{theorem}

\begin{theorem}[непрерывность сложной функции]
Пусть
\[
\lim_{x\to\alpha} f(x)=\beta
\]
и функция \(g\) непрерывна в точке \(\beta\). Тогда
\[
\lim_{x\to\alpha} g(f(x))=g(\beta).
\]
\end{theorem}

\begin{proof}
Возьмём произвольную последовательность \(\{x_n\}\) такую, что
\[
x_n\to\alpha,\qquad x_n\ne\alpha.
\]
По определению предела по Гейне
\[
f(x_n)\to\beta.
\]
Так как \(g\) непрерывна в точке \(\beta\), по определению непрерывности
\[
g(f(x_n))\to g(\beta).
\]
Следовательно, по определению предела по Гейне
\[
\lim_{x\to\alpha} g(f(x))=g(\beta).
\]
\end{proof}

\begin{corollary}
Если \(y=f(x)\) непрерывна в точке \(\alpha\), \(z=g(y)\) непрерывна
в точке \(\beta\), причём \(\beta=f(\alpha)\), то сложная функция
\(h(x)=g(f(x))\) непрерывна в точке \(\alpha\), и
\[
\lim_{x\to\alpha} h(x)
=
\lim_{x\to\alpha} g(f(x))
=
g(f(\alpha))
=
h(\alpha).
\]
\end{corollary}

\begin{theorem}[устойчивость знака]
Если функция \(f\) непрерывна в точке \(\alpha\) и \(f(\alpha)\ne 0\), то
существует \(\delta>0\) такое, что для всех \(x\in(\alpha-\delta,\alpha+\delta)\)
выполнено
\[
f(x)\cdot f(\alpha)>0.
\]
\end{theorem}

\begin{proof}
Так как \(f(\alpha)\ne 0\), возьмём
\[
\varepsilon=\frac{|f(\alpha)|}{2}>0.
\]
По непрерывности \(f\) в точке \(\alpha\) найдётся \(\delta>0\) такое, что
для всех \(x\) с условием \(|x-\alpha|<\delta\) выполнено
\[
|f(x)-f(\alpha)|<\varepsilon.
\]
Тогда
\[
f(x)\in\bigl(f(\alpha)-\varepsilon,\ f(\alpha)+\varepsilon\bigr),
\]
причём этот интервал не содержит нуля и целиком лежит по ту же сторону
от нуля, что и \(f(\alpha)\). Следовательно, \(f(x)\) имеет тот же знак,
что и \(f(\alpha)\), то есть
\[
f(x)\cdot f(\alpha)>0.
\]
\end{proof}

\Section{Ограниченность функции, непрерывной на отрезке}

\begin{theorem}[Первая теорема Вейерштрасса]
Если функция \(f(x)\) непрерывна на отрезке \([a,b]\), то она ограничена на этом
отрезке:
\[
[f(x)\in C[a,b]]\ \Longrightarrow\ [f(x)\text{ ограничена на }[a,b]].
\]
\end{theorem}

\begin{proof}
Предположим противное: пусть \(f(x)\) не ограничена на \([a,b]\). Тогда
\[
\forall n\in\mathbb N\ \exists x_n\in[a,b]:\quad |f(x_n)|>n.
\]
Последовательность \(\{x_n\}\) ограничена, так как все её члены лежат на отрезке
\([a,b]\). По теореме Больцано--Вейерштрасса из неё можно выделить сходящуюся
подпоследовательность:
\[
\exists \{x_{n_k}\}:\quad x_{n_k}\to c\in[a,b].
\]
Так как \(f\) непрерывна на \([a,b]\), в частности в точке \(c\), то
\[
\lim_{k\to\infty} f(x_{n_k})=f(c).
\]
Следовательно, последовательность \(\{f(x_{n_k})\}\) сходится, а значит,
ограничена. Но с другой стороны,
\[
|f(x_{n_k})|>n_k\to\infty,
\]
то есть \(\{f(x_{n_k})\}\) неограничена. Полученное противоречие показывает, что
предположение неверно. Значит, \(f(x)\) ограничена на \([a,b]\).
\end{proof}

\Section{Достижение точной верхней и точной нижней граней функции,
непрерывной на отрезке}
\begin{theorem}[Вторая теорема Вейерштрасса]
Если функция \(f(x)\) непрерывна на отрезке \([a,b]\), то она достигает на этом
отрезке своей точной верхней и точной нижней граней:
\[
[f(x)\in C[a,b]]
\ \Longrightarrow\
\bigl[\exists \alpha,\beta\in[a,b]:\
\sup_{[a,b]} f(x)=f(\alpha),\quad
\inf_{[a,b]} f(x)=f(\beta)\bigr].
\]
\end{theorem}

\begin{proof}
Обозначим
\[
M=\sup_{[a,b]} f(x).
\]
Предположим противное: пусть значение \(M\) не достигается, то есть
\[
f(x)<M\quad \forall x\in[a,b].
\]
Тогда функция
\[
g(x)=\frac{1}{M-f(x)}
\]
непрерывна на \([a,b]\), так как знаменатель не обращается в нуль. По первой
теореме Вейерштрасса \(g(x)\) ограничена на \([a,b]\), то есть
\[
\exists A>0:\quad g(x)\le A\quad \forall x\in[a,b].
\]
Следовательно,
\[
\frac{1}{M-f(x)}\le A
\ \Longrightarrow\
M-f(x)\ge \frac{1}{A}
\ \Longrightarrow\
f(x)\le M-\frac{1}{A}<M.
\]
Но это означает, что число \(M-\frac{1}{A}\) является верхней гранью \(f(x)\) на
\([a,b]\), меньшей чем \(M\), что противоречит определению \(M\) как точной
верхней грани.

Полученное противоречие показывает, что существует точка \(\alpha\in[a,b]\), в
которой
\[
f(\alpha)=M=\sup_{[a,b]} f(x).
\]

Аналогично доказывается существование точки \(\beta\in[a,b]\), в которой
достигается точная нижняя грань:
\[
f(\beta)=\inf_{[a,b]} f(x).
\]
Теорема доказана.
\end{proof}

\Section{Теорема о промежуточных значениях непрерывной функции}


\begin{theorem}[Больцано--Коши]
Если функция \(f(x)\) непрерывна на отрезке \([a,b]\) и
\[
f(a)\cdot f(b)<0,
\]
то существует точка \(c\in(a,b)\) такая, что
\[
f(c)=0.
\]
\end{theorem}

\begin{proof}
Пусть для определённости \(f(a)<0\), \(f(b)>0\). Рассмотрим множество
\[
X=\{x\in[a,b]\mid f(x)<0\}.
\]
Оно непусто, так как \(a\in X\), и ограничено сверху, поэтому существует
\[
c=\sup X.
\]
Покажем, что \(f(c)=0\) и \(c\in(a,b)\).

Так как \(f(a)<0\) и функция \(f\) непрерывна в точке \(a\), существует
\(\delta_1>0\) такое, что для всех \(x\in(a,a+\delta_1)\) выполнено
\(f(x)<0\). Значит, \(c>a\).

Аналогично, так как \(f(b)>0\), существует \(\delta_2>0\) такое, что для всех
\(x\in(b-\delta_2,b)\) выполнено \(f(x)>0\). Значит, \(c<b\).

Следовательно, \(c\in(a,b)\).

Предположим, что \(f(c)\ne 0\). Тогда либо \(f(c)<0\), либо \(f(c)>0\).

Если \(f(c)<0\), то по непрерывности \(f\) в точке \(c\) найдётся окрестность
\((c-\delta,c+\delta)\), в которой \(f(x)<0\). Тогда точки, большие \(c\),
принадлежали бы множеству \(X\), что противоречит тому, что
\(c=\sup X\).

Если \(f(c)>0\), то найдётся окрестность \((c-\delta,c+\delta)\), в которой
\(f(x)>0\). Тогда все точки, достаточно близкие к \(c\) слева, не принадлежали
бы \(X\), что противоречит определению \(c\) как точной верхней грани.

Полученное противоречие показывает, что \(f(c)=0\).
\end{proof}

\begin{theorem}[О промежуточных значениях]
Пусть функция \(y=f(x)\) непрерывна на отрезке \([a,b]\), и пусть
\[
\alpha=f(a),\qquad \beta=f(b).
\]
Тогда для любого числа \(\gamma\), лежащего между \(\alpha\) и \(\beta\)
(то есть \(\gamma\in[\alpha,\beta]\)), найдётся точка \(c\in(a,b)\) такая, что
\[
f(c)=\gamma.
\]
\end{theorem}

\begin{proof}
Рассматриваем вспомогательную функцию
\[
g(x)=f(x)-\gamma.
\]
Она непрерывна на \([a,b]\), а на концах отрезка принимает значения разных знаков:
\[
g(a)=f(a)-\gamma=\alpha-\gamma,\qquad
g(b)=f(b)-\gamma=\beta-\gamma.
\]
По теореме Больцано--Коши о корне непрерывной функции существует точка
\(c\in(a,b)\), в которой \(g(c)=0\), то есть \(f(c)=\gamma\).
\end{proof}